\documentclass[11pt,fleqn]{report}
\usepackage{graphicx}
\usepackage{amssymb}
\usepackage{longtable}
\usepackage[T1]{fontenc}
\usepackage{lmodern}
\usepackage{cml}
\lstnewenvironment{codeXML}
    {\lstset{basicstyle=\small\ttfamily,
             language=XML,
             xleftmargin=0.0in,
             xrightmargin=0.0in,
             backgroundcolor=\color{Orchid!60}, % 60% orchid, 40% white
             escapeinside={\%*}{*)}, %  .tex file characters inside %*...*) are
                                     %  treated as Tex commands.
            }
    }
    { }

\newcommand{\ModelName}{CML Message Publisher}
\newcommand{\ModelAuthor}{Gary Turner}
\newcommand{\ModelAffiliation}{Odyssey Space Research}
\newcommand{\ModelDate}{March 2023}
\pagestyle{plain}
\usepackage[parfill]{parskip}
\usepackage[margin=1.5cm]{caption}
\usepackage{subfig}

\begin{document}
\titlePage
\clearpage % Reset page number
\pagenumbering{roman}
\tableofcontents % Print table of contents
\vfill
{\Large \textbf {Revision History}}
{\large
\begin{center}
\begin{tabular}{|l||l|l|l|}
\hline
 \textbf{Version} & \textbf{Date} & \textbf{Author} & \textbf{Purpose} \\ \hline
 \hline
1 & March 2023 & Gary Turner & Initial Version \\ \hline
 \hline
2 & August 2025 & Alexandre Masset & V\&V and Adding Status to Tests and Docs \\ \hline
 \hline
3 & April 2026 & Nino Tarantino & Refactored for testing \\ \hline
\end{tabular}
\end{center}
}

\chapter{Introduction} % Make a "chapter" heading
\pagenumbering{arabic} % Start text with arabic 1
The CML Message Publishing model is an upgrade from the JEOD MessageHandler
model. Like the JEOD version, this implementation includes a set of tunable
message levels:
\begin{itemize}
\item Fail represents a critical failure. The sim switches off. This level
cannot be ignored.
\item Error represents the case where a significant error has been detected.
This is not terminal, the sim can continue to run without seg-faulting, but all
data beyond the publishing of this error message is suspect.
\item Warning respresents the case where a data anomaly has been detected.
Typically, the code has logic to bypass the problem but users should be aware
of the potential for data invalidation, and should fix the situation /
configuration.
\item Status represents a message sent to the screen to inform the user of some activity 
with abbreviated logs. The line and file are not provided in this case.
\item Inform represents a message sent to the screen to inform the user of some
activity within the simulation. This does not mean there has been an adverse
impact on the simulation data, it is for informational purposes only.
\item Debug are comments included to assist with debugging only. This level is
rarely used.
\end{itemize}
The model can be tuned so that only messages more severe than some level
actually get published.

\chapter {Requirements}
The CML Message Publishing model has the following requirements:
\begin{enumerate}
\item Reports shall be generated with sufficient speed to avoid
      adversely affecting simulation performance or real time operations. 
\item Reports shall include: 
  \begin{enumerate}
  \item a message significance/priority classification
        (e.g. critical fault, error, warning, etc.); 
  \item a timestamp indicating a reproduceable clock value
        (e.g. Dynamic-time) at which they were generated; 
  \item location information on the source of the message,
        including the file name and line number; 
  \end{enumerate}
\item The simulation shall provide a mechanism to block messages below a
      user-settable priority threshold.
\end{enumerate}



\chapter {Model Specification}
\section{Message Importance}
The model contains an enumeration of the 6 levels of message importance. This
can be extended if more levels are deemed necessary.
The model includes an instance of this enumeration and a setter method to allow
the user to show (or hide) messages more (or less) significant than some
desired threshold level of importance.

\section{Message Publishing}
The model provides a single method using a variadic template to publish a
message: \textit{CMLMessage::publish(...)}. It takes the following arguments:
\begin{itemize}
\item The importance of the message being sent
\item The filename from which the message was generated
\item The line number within the file where the message was generated
\item Any number of output-streamable arguments of any data type for which \textit{std::ostringstream::operator<<} is defined
\end{itemize}

The model also provides interface methods that effectively alias this 
common method and support omission of the first argument:
\begin{itemize}
\item \textit{CMLMessage::fail(...)} calls
      \textit{CMLMessage::publish( CMLMessage::Fail, ...)}.
\item \textit{CMLMessage::error(...)} calls
      \textit{CMLMessage::publish( CMLMessage::Error, ...)}.
\item \textit{CMLMessage::warn(...)} calls
      \textit{CMLMessage::publish( CMLMessage::Warning, ...)}.
\item \textit{CMLMessage::inform(...)} calls
      \textit{CMLMessage::publish( CMLMessage::Inform, ...)}.
\item \textit{CMLMessage::debug(...)} calls
      \textit{CMLMessage::publish( CMLMessage::Debug, ...)}.
\item \textit{CMLMessage::status(...)} calls
      \textit{CMLMessage::publish( CMLMessage::Status, ...)}
\end{itemize}
These aliases appear very similar to the method calls supported by JEOD's
MessageHandler class. Thus transitioning from existing JEOD messaging to
CML messaging should be possible by changing the header file include
statement, and replacing \textit{jeod::MessageHandler} with \textit{CMLMessage}.

This model also provides a helper method \textit{CMLMessage::printf\_fmt} to print
outputs formatted like a printf statement. This method is not associated with any specific
message type.

\section {Termination}
The method \textit{terminate(...)} is called internally when a message is
published with Fail (critical fault) level. For Trick-based simulations, this
results in a shutdown of the simulation. When including the model from a non-Trick
context, an additional message is printed informing the user that a shutdown would have
occurred, but execution will continue.

\chapter {User's Guide}
\section {Implementation}
The model's contents are free functions placed within the \textit{CMLMessage}
namespace. Each of these functions calls a non-variadic version of
\textit{CMLMessage::publish(...)} under the hood. An instantiation of the
\textit{CMLMessage::PublishLevel} enumeration is contained within an anonymous
namespace within \textit{src/cml\_message.cc} which governs the minimum criticality
of messages that will be published and can be modified by calling the
\textit{CMLMessage::set\_publish\_level(...)} function.

\section{Initialization}
The model has no initialization requirement.

\section {Routine Execution}
The model's \textit{publish(...)} method gets called from a model upon
detection of some feature or anomaly worth reporting. Typically, the arguments
are as outlined above (and repeated here):
\begin{itemize}
\item The importance of the message being sent
\item The filename from which the message was generated
\item The line number within the file where the message was generated
\item Any number of output-streamable arguments of any data type for which \textit{std::ostringstream::operator<<} is defined
\end{itemize}

For example:
\begin{codeC}
  CMLMessage::publish( CMLMessage::Error,
                       __FILE__,
                       __LINE__,
                       "Some message about variable ", var.name,
                       " whose value ", var.value, " is larger than some "
                       "threshold, ", var.threshold);
\end{codeC}

Alternatively, the interface methods can be used; the following call has the
same effect:
\begin{codeC}
  CMLMessage::error( __FILE__,
                     __LINE__,
                     "Some message about variable ", var.name,
                     " whose value ", var.value, " is larger than some "
                     "threshold, ", var.threshold);
\end{codeC}

\textbf{Notes:}
\begin{itemize}
\item In the arguments, a single string argument may be spaced across
multiple lines without needing comma separation. In the examples above, the
string \textit{" is larger than some threshold, "} is split across two lines.

\item Spacing between arguments in the final collation of arguments must be
manually provided. In the examples above, the string \textit{" whose value "}
starts and ends with a space to create a separation between it and the
values \textit{var.name} before it, and \textit{var.value} after it.
\end{itemize}

The verification simulation includes the following code:
\begin{codeC}
    int ii = 4;
    std::string part2 = "is a collection of ";
    CMLMessage::publish(
        CMLMessage::Error,
        __FILE__, __LINE__,
        "This ",part2,ii," arguments");
\end{codeC}

\begin{figure}[h]
  \centering
  \includegraphics[scale=0.6]{graphics/error_message.png}
  \caption{Example Error Message}
  \label{fig:example_message}
\end{figure}

The \textit{CMLMessage::printf\_fmt(std::string format, Args... args)} can be used
to print a formatted string. Typical use is as follows:
\begin{codeC}
      CMLMessage::printf_fmt("%12.6g  %14.12g  %18.12E  %15.3f   %6d", 
                                    x,       x,       x,      x,     ii);
\end{codeC}

This method can be combined with the \textit{publish(...)} method as in:
\begin{codeC}
      CMLMessage::debug(
          __FILE__, __LINE__,
          "Testing printf_fmt: ",
          CMLMessage::printf_fmt("%12.6g  %14.12g  %18.12E  %15.3f   %6d",
                                    x,       x,       x,      x,     ii),
          "\nUnformatted value of x:",x);
\end{codeC}

\newpage

\section {Optional Execution}
The threshold level for message publishing can be set using
the \textit{set\_publish\_level( PublishLevel)} method
\begin{codePython}
trick.CMLMessage.set_publish_level (trick.CMLMessage.Warning)
\end{codePython}

\section {Extension}
The model provides the same five severity-level categories traditionally used
by the JEOD message-handler system. A new message category \textbf{Status} was
also added. These categories can be further extended if the user has write-acess
to the source code.

Suppose we want to introduce a new message category, let's call it
\textbf{Alert}. This can be acomplished with the following steps;
\begin{itemize}
\item Decide the hierarchy of where \textbf{Alert} fits with the other
categories, e.g. between \textit{Error} and \textit{Warning}.
\item Insert the new category into the \textit{PublishLevel} enumeration.
Recall that the hierarchy is used to identify which messages
should be published to standard out so it is essential that the new category
has appropriate placing. If it is simply (and possibly erroneously) placed
at the end, then setting the publish-level to \textit{Debug} would block
messages identified with \textbf{Alert}.
\begin{codeC}
  enum PublishLevel {
    Fail = 0,
    Error,
    Alert,
    Warning,
    Status,
    Inform,
    Debug
  };
\end{codeC}
\item (optional) provide an interface method if desired.
\begin{codeC}
  template<typename... Args>
  static void alert( const std::string& file,
                     int                line,
                     Args&&...          args)
  {
    const std::string text = collate_args(std::forward<Args>(args)...);
    publish(PublishLevel::Alert, file, line, text);
  }
\end{codeC}
\end{itemize}

\chapter{Verification}
\section{Code Coverage}

\begin{figure}[h]
  \centering
  \includegraphics[width=\textwidth]{graphics/coverage.png}
  \caption{Code Coverage}
  \label{fig:Code Coverage}
\end{figure}

\subsection{Exceptions}

\begin{itemize}
      \item \textbf{Lines 33-35}: Unreachable in a Trick simulation environment because the simulation has already terminated, and unreachable in a unit test environment because the termination call is mocked out. See Figure \ref{fig:missing_coverage}.
\end{itemize}

\begin{figure}[h]
  \centering
  \includegraphics[width=0.7\textwidth]{graphics/uncovered.png}
  \caption{Missing Source File Line Coverage}
  \label{fig:missing_coverage}
\end{figure}

\section{Unit Tests}
The GoogleTest testing and mocking framework (GTest) is used to test the main
functionality of the model. The following tests are defined:
\begin{itemize}
\item \textbf{SetPublishLevel}: Ensures that the proper publish level is set when the user calls the \textit{CMLMessage::set\_publish\_level(...)} method.
\item \textbf{CollateArgs}: Tests the aspect of the model which combines output-streamable arguments into a single \textit{std::string}.
\item \textbf{InterfaceMethods}: Ensures that the proper invocation of \textit{CMLMessage::publish(...)} is dispatched when calling an interface method such as \textit{CMLMessage::error(...)}.
\item \textbf{PrintFormat}: Tests the print formatting functionality of the model.
\end{itemize}

\section{SIM\_verif}
The model also includes a simple verification simulation.
The verification simulation has two runs; there is no verification data
because the model's purpose is entirely to publish data to standard-out.

\subsection{Model Capability}

The verification simulation has one main run which allows for a visual inspection of the message
outputs at each publishing level.

For the main run, \textit{RUN\_test}, on each of the 5 timesteps a "sample model" sends the
same 5 messages to the Message Publisher, one message for each category except for critical failure (Fail). 

The publishing threshold is gradually made more inclusive:
\begin{itemize}
\item t=0 Only Fail messages are published; no output.
\item t=1 The threshold is set to Error. The Error message is published and the
rest ignored.
\item t=2 The threshold is set to Warning. The Error and Warning messages are
published and the rest ignored.
\item t=3 The threshold is set to Status. The Error, Warning and Status messages
are published and the rest ignored.
\item t=4 The threshold is set to Inform. The Error, Warning, Inform and Status messages
are published and only the Debug message ignored.
\item t=5 The threshold is set to Debug. All 5 messages are published.
\end{itemize}

Also at t=5, a message is sent with level Fail. the simulation is terminated
by this message.

The last round of outputs is shown in Figure \ref{fig:verif_output}.
\begin{figure}[h]
  \centering
  \includegraphics[width=\textwidth]{graphics/verif_sim_output.png}
  \caption{Sim Output}
  \label{fig:verif_output}
\end{figure}

\clearpage

\subsection{Model speed}
One of the requirements addresses the need for the message publisher to
process sufficiently quickly to avoid impacting simulation speed.

The run 
\textit{RUN\_speed\_test} simply processes through a for loop, sending an
error message 10,000 times. This processes all messages in 
approximately 3 seconds per the linux \textit{time} utility. This
satisfactorily meets the requirement; it is not expected that any simulation
would be producing that many messages.

For comparison, the same exercise was processed with the JEOD message
handler. It runs marginally faster (2.5 seconds for the same 10,000
messages). With the number of messages reduced to 1,000, the JEOD Message
Handler runs a little slower (0.5 seconds) than this model (0.4 seconds).
\end{document}
